\chapter{Wilson Lattice Gauge Theory}
\label{ch:constructive-wilson-lattice-gauge-theory}

\ConventionsEuclidean

Wilson lattice gauge theory gives a gauge-invariant finite-regulator
definition of compact gauge fields.  Its continuum use requires a scaling
trajectory, reflection positivity, locality, and a reconstruction or
continuum-limit argument.  The lattice theory is not merely a discretized
notation for continuum perturbation theory; it is a nonperturbative
regulated path integral.

\section{Link Variables and Gauge Transformations}

Let \(\Lambda\subset a\mathbb Z^D\) be a finite hypercubic lattice and \(G\)
a compact Lie group.  To each oriented link \(\ell=(x,x+a\hat\mu)\), assign
\[
  U_\mu(x)\in G,
\]
with \(U_{-\mu}(x+a\hat\mu)=U_\mu(x)^{-1}\).  A lattice gauge transformation
is a function \(g:\Lambda^0\to G\) acting by
\[
  U_\mu(x)\mapsto g(x)U_\mu(x)g(x+a\hat\mu)^{-1}.
\]
The plaquette holonomy is
\[
  U_{\mu\nu}(x)
  =
  U_\mu(x)U_\nu(x+a\hat\mu)
  U_\mu(x+a\hat\nu)^{-1}U_\nu(x)^{-1}.
\]

\section{Wilson Action}

Fix a faithful unitary representation \(R\) and normalize the trace in the
continuum convention of the gauge-theory volume.  A Wilson plaquette action
has the form
\[
  S_W[U]
  =
  \beta\sum_{p}
  \left(
  1-\frac{1}{\dim R}\operatorname{Re}\operatorname{Tr}_R U_p
  \right).
\]
The finite-volume partition function is the compact group integral
\[
  Z=\int\prod_{\ell}dU_\ell\,e^{-S_W[U]},
\]
where \(dU\) is Haar measure.

\section{A Finite \texorpdfstring{\(\mathbb Z_2\)}{Z2} Gauge Benchmark}

Before taking any continuum limit, a lattice gauge theory is an ordinary
finite-dimensional integral.  The smallest nontrivial pure gauge benchmark is
the \(\mathbb Z_2\) gauge model on a finite periodic cubic lattice.  Let
\(E\) be the set of oriented positive links, let \(P\) be the set of
elementary plaquettes, and assign
\[
  U_e\in\{\pm1\},\qquad e\in E .
\]
For a plaquette \(p\), define
\[
  U_p=\prod_{e\in\partial p}U_e ,
\]
where orientation is irrelevant because \(U_e^{-1}=U_e\).  The finite
partition function at parameter \(\beta\in\mathbb R\) is
\[
  Z_\Lambda(\beta)
  =
  \sum_{U:E\to\{\pm1\}}
  \exp\!\left(\beta\sum_{p\in P}U_p\right).
\]
A gauge transformation is a function \(g:V\to\{\pm1\}\) acting on a link
\(e=(x,y)\) by
\[
  U_e\longmapsto g_xU_eg_y .
\]
The plaquette variables and Wilson loop variables
\[
  W(C)=\prod_{e\in C}U_e
\]
for closed lattice contours \(C\) are gauge invariant.

\begin{proposition}[Finite \(\mathbb Z_2\) character expansion]
\label{prop:z2-gauge-character-expansion}
Let \(t=\tanh\beta\).  For a closed contour \(C\), regard \(C\) as a
\(\mathbb Z_2\)-one-chain, and regard a subset \(A\subset P\) as a
\(\mathbb Z_2\)-two-chain with boundary \(\partial A\).  Then
\[
  Z_\Lambda(\beta)
  =
  2^{|E|}(\cosh\beta)^{|P|}
  \sum_{\substack{A\subset P\\ \partial A=0}} t^{|A|},
\]
and
\[
  \langle W(C)\rangle_{\Lambda,\beta}
  =
  \frac{\displaystyle
  \sum_{\substack{A\subset P\\ \partial A=C}}t^{|A|}}
  {\displaystyle
  \sum_{\substack{A\subset P\\ \partial A=0}}t^{|A|}} .
\]
The sums are finite sums over plaquette subsets.  Thus the strong-coupling
area mechanism in this model is not an analogy: at finite cutoff the Wilson
loop numerator is exactly a weighted sum over lattice surfaces whose
boundary is the loop.
\end{proposition}

\begin{proof}
For each plaquette variable \(U_p=\pm1\),
\[
  e^{\beta U_p}=\cosh\beta\,(1+tU_p)
  =
  \cosh\beta\sum_{n_p=0}^1 t^{n_p}U_p^{n_p}.
\]
Multiplying over plaquettes gives
\[
  \exp\!\left(\beta\sum_p U_p\right)
  =
  (\cosh\beta)^{|P|}
  \sum_{A\subset P}
  t^{|A|}
  \prod_{p\in A}\prod_{e\in\partial p}U_e .
\]
Summing over a single link variable gives
\(\sum_{U_e=\pm1}U_e^m=0\) if \(m\) is odd and \(2\) if \(m\) is even.
Therefore the partition function keeps precisely those plaquette subsets for
which every link appears an even number of times, namely
\(\partial A=0\), and contributes the overall factor \(2^{|E|}\).

For the Wilson-loop numerator, the additional factor
\(\prod_{e\in C}U_e\) changes the parity condition on each link from
\((\partial A)_e=0\) to \((\partial A)_e=C_e\).  Dividing the resulting
finite numerator by the finite partition function gives the displayed
expectation value.
\end{proof}

The companion script
\[
  \texttt{qft\_scripts/z2\_gauge\_3d\_metropolis.py}
\]
samples this finite model by single-link Metropolis updates and measures the
plaquette average and rectangular Wilson loops.  Its output is finite-cutoff
data.  An infinite-volume transition, a continuum limit, or a continuum
confinement statement requires additional limiting hypotheses and estimates.

\section{Surface Counting and the Strong-Coupling Area Bound}

The exact formula in
Proposition~\ref{prop:z2-gauge-character-expansion} becomes an area-law
estimate only after one controls how many surfaces of a given area have the
same boundary.  This counting step is often hidden in the phrase
``strong-coupling expansion.''  We spell out the finite statement.

Fix a finite lattice and a closed contour \(C\) which is a boundary in the
\(\mathbb Z_2\) plaquette chain complex.  Define the lattice area
\[
  A_{\min}(C)
  =
  \min\{|A|:A\subset P,\ \partial A=C\}.
  \label{eq:z2-minimal-lattice-area}
\]
For \(n\geq0\), let
\[
  N_C(n)
  =
  \#\{A\subset P:\partial A=C,\ |A|=A_{\min}(C)+n\}.
  \label{eq:z2-surface-counting-number}
\]
Similarly let \(N_0(n)\) count closed plaquette subsets of area \(n\).

\begin{proposition}[Finite surface expansion and area estimate]
\label{prop:z2-surface-counting-area-bound}
Let \(\beta\geq0\) and \(t=\tanh\beta\).  If \(C\) is a boundary, then
\[
  \langle W(C)\rangle_{\Lambda,\beta}
  =
  t^{A_{\min}(C)}
  \frac{\sum_{n\geq0}N_C(n)t^n}
       {\sum_{n\geq0}N_0(n)t^n}.
  \label{eq:z2-surface-counting-expansion}
\]
In particular, since the denominator is at least \(1\),
\[
  0\leq
  \langle W(C)\rangle_{\Lambda,\beta}
  \leq
  t^{A_{\min}(C)}
  \sum_{n\geq0}N_C(n)t^n .
  \label{eq:z2-area-bound-before-entropy}
\]
If, for a family of contours, the surface entropy obeys
\[
  N_C(n)\leq K(C)\rho^n
  \qquad(n\geq0)
  \label{eq:z2-surface-entropy-hypothesis}
\]
with \(0\leq\rho t<1\), then
\[
  \langle W(C)\rangle_{\Lambda,\beta}
  \leq
  \frac{K(C)}{1-\rho t}\,
  \exp\{-[-\log t]A_{\min}(C)\}.
  \label{eq:z2-entropy-area-bound}
\]
\end{proposition}

\begin{proof}
The numerator in
Proposition~\ref{prop:z2-gauge-character-expansion} is a finite sum over
plaquette subsets satisfying \(\partial A=C\).  Grouping these subsets by
the excess area \(n=|A|-A_{\min}(C)\) gives the numerator
\[
  t^{A_{\min}(C)}\sum_{n\geq0}N_C(n)t^n .
  \]
The same grouping with \(C=0\) gives the denominator.  For
\(\beta\geq0\), all terms in the denominator are nonnegative and the empty
closed surface gives the term \(1\); hence the denominator is at least
\(1\), proving \eqref{eq:z2-area-bound-before-entropy}.  The final estimate
is the geometric-series bound
\[
  \sum_{n\geq0}N_C(n)t^n
  \leq
  K(C)\sum_{n\geq0}(\rho t)^n
  =
  \frac{K(C)}{1-\rho t}.
  \]
\end{proof}

For a rectangular loop \(C_{R,T}\) in a coordinate plane,
\(A_{\min}(C_{R,T})=RT/a^2\) when the rectangle is smaller than the chosen
finite box and the flat spanning rectangle is the unique minimal surface in
that homology class.  Equation~\eqref{eq:z2-entropy-area-bound} then gives a
finite-regulator area bound at sufficiently small \(t\).  If the constants
in the surface-counting estimate have at most perimeter growth,
\(K(C_{R,T})\leq \exp\{c(R+T)/a\}\), the area coefficient in lattice units is
bounded below by
\[
  \sigma_{\rm lat}(\beta)\geq -\log(\tanh\beta)
\]
with the displayed perimeter and finite entropy prefactors kept separately.
This is a strong-coupling theorem at fixed cutoff.  It is not a proof of
continuum confinement, because a continuum Yang--Mills limit in four
dimensions would require control along a scaling trajectory where the Wilson
parameter does not remain in this small-\(\beta\) domain.

\begin{example}[One-cube surface polynomial]
\label{ex:z2-one-cube-surface-polynomial}
On the boundary of one elementary cube, let \(C\) be the boundary of one
face.  There are exactly two plaquette subsets with boundary \(C\): the face
itself, with area \(1\), and the complementary five faces, with area \(5\).
The closed subsets are the empty surface and the six-face cube boundary.
Therefore
\[
  \langle W(C)\rangle
  =
  \frac{t+t^5}{1+t^6}.
  \label{eq:z2-one-cube-wilson-polynomial}
\]
This example shows explicitly why the leading strong-coupling term is the
minimal spanning surface and why closed surfaces correct both numerator and
denominator.
\end{example}

\section{Continuum Expansion}

For a smooth continuum connection \(A_\mu(x)\), set
\[
  U_\mu(x)=\exp(a A_\mu(x)).
\]
The plaquette has expansion
\[
  U_{\mu\nu}(x)
  =
  \exp\!\left(a^2F_{\mu\nu}(x)+O(a^3)\right).
\]
If the representation trace obeys
\[
  \operatorname{Tr}_R(XY)=\kappa_R\,\operatorname{tr}(XY),
\]
then
\[
  1-\frac{1}{\dim R}\operatorname{Re}\operatorname{Tr}_R U_{\mu\nu}
  =
  -\frac{a^4\kappa_R}{2\dim R}
  \operatorname{tr}(F_{\mu\nu}F_{\mu\nu})
  +O(a^6)
\]
for anti-Hermitian \(A_\mu\).  Matching to
\[
  \frac{1}{4g^2}\int d^Dx\,\operatorname{tr}(F_{\mu\nu}F_{\mu\nu})
\]
fixes the relation between \(\beta\), \(g\), and the trace normalization.

\section{Wilson Loops}

For a closed oriented lattice contour \(C\),
\[
  W_R(C)=\operatorname{Tr}_R\prod_{\ell\in C}U_\ell
\]
is gauge invariant.  Wilson loops are the basic gauge-invariant probes of
line charges.  Their strong-coupling character expansion gives the area-law
mechanism described in Chapter~\ref{ch:global-confinement-screening-oblique}.
Their continuum limit requires line renormalization and an operator scheme.

\section{Reflection Positivity and Transfer Matrix}

For the Wilson plaquette action, reflection positivity follows by assigning
plaquettes to the two halves of the lattice and expanding cross-plane
plaquette Boltzmann weights in positive characters.  The reflection-positive
inner product on positive-time functions gives a transfer matrix
\[
  T=e^{-aH}
\]
with \(H\geq0\) after quotienting null vectors.  Gauge constraints are
implemented by the gauge-invariant subspace or by a gauge-covariant transfer
matrix followed by projection.

\section{Fermions and Chiral Difficulties}

Adding lattice fermions requires a fermion action whose reflection
positivity, locality, spectrum, and continuum chiral properties are checked
separately.  The mid-link fermion reflection structure of
Chapter~\ref{ch:lattice-reflection-positivity} is one such finite-regulator
positivity mechanism.  Chiral gauge theories require more: gauge invariance,
locality, anomaly cancellation, and a definition of the fermion determinant
phase compatible with the target continuum theory.

\begin{openproblem}[Lattice continuum QCD construction]
\label{op:lattice-continuum-qcd-construction}
Construct four-dimensional continuum Yang--Mills or QCD from Wilson-type
lattice gauge theory with a mass gap, OS reconstruction, local gauge-invariant
operator limits, and controlled line-operator and theta-angle sectors.
\end{openproblem}
